\subsection{Einbettung des GNN als MILP}
\label{sec:gnn-milp-einbettung}

Nach dem Training werden die Gewichte und Biasparameter des GNN fixiert und
das resultierende Netz als gemischt-ganzzahliges lineares Modell in das
FJSP-Grundmodell eingebettet. Dadurch gehen die vom GNN geschätzten
Pünktlichkeitswahrscheinlichkeiten unmittelbar in die Optimierung ein. Die
Einbettung folgt dabei dem üblichen Vorgehen, affine Netzwerkschichten direkt
als lineare Gleichungen abzubilden und die ReLU-Aktivierungen mithilfe von
Binärvariablen exakt zu linearisieren \cite{anderson2020strong}.

Untersucht werden drei Architekturvarianten mit identischen Knotenfeatures,
Dimensionen der verdeckten Repräsentationen, Aktivierungsfunktionen sowie
Pooling- und Ausgabeschichten. Sie unterscheiden sich ausschließlich durch
die im Message Passing verwendete Graphstruktur. Die lineare Variante (LIN)
verarbeitet nur den Zustand des jeweiligen Knotens. Die JOB-Variante ergänzt feste
Jobpräzedenzkanten. Die FULL-Variante berücksichtigt darüber hinaus
Maschinenkanten, die erst durch die Maschinenzuordnungs- und
Reihenfolgeentscheidungen des FJSP bestimmt werden. Die technologisch
vorgegebenen Jobpräzedenzkanten bleiben auch in der FULL-Variante fest; nur die
Maschinenkanten sind entscheidungsabhängig.

Während des Trainings liegt die Graphstruktur jedes Trainingsbeispiels nach
Festlegung des zugehörigen Ablaufplans vollständig vor. Bei der Einbettung in
das Optimierungsmodell gilt dies lediglich für die Jobpräzedenzkanten. In der
FULL-Variante wird die aktive Maschinennachbarschaft dagegen endogen während
der Lösung bestimmt. Deshalb erfordert nur diese Variante zusätzliche
Hilfsvariablen zur Kopplung der binären Kantenentscheidungen mit den
kontinuierlichen Knotenrepräsentationen. Die Architektur der beiden
Message-Passing-Varianten orientiert sich an GraphSAGE
\cite{hamilton2017inductive}; anstelle einer Mittelwertbildung wird hier eine
Summenaggregation verwendet.

\begin{table}[H]
\centering
\caption{Ergänzende Notation der GNN-Formulierung}
\label{tab:notation_fjsp_gnn}
\scriptsize
\setlength{\tabcolsep}{3pt}
\renewcommand{\arraystretch}{0.88}
\begin{tabularx}{\textwidth}{@{}lX@{}}
\toprule
\multicolumn{2}{l}{\textbf{Mengen und Parameter}}\\
\midrule

$\mathcal{A}=\{\mathrm{LIN},\mathrm{JOB},\mathrm{FULL}\}$
& GNN-Architekturvarianten \\

$\mathcal{L}=\{1,\dots,L\}$
& GNN-Schichten; hier: $L=1$ \\

$\mathcal{Q}_0=\{1,\dots,5\}$
& Eingangskanäle bzw. Knotenfeatures \\

$\mathcal{Q}_{\ell}=\{1,\dots,d_{\ell}\}$
& Hidden-Kanäle der Schicht $\ell$; hier: $d_{\ell}=4$ \\

$\mathcal{P}_i$
& direkte Jobvorgänger von Operation $i$ \\

$\mathcal{O}_u$
& Operationen des Jobs $u$ \\

$\mathcal{E}^{M}$
& zulässige Maschinenkanten $(j,i,k)$ \\

$\alpha_u$
& gefordertes Serviceniveau für Job $u$ \\

$W_{\mathrm{lin}}^{(\ell,\mathrm{LIN})}$
& Gewichtsmatrix der linearen Schicht $\ell$ \\

$W_{\mathrm{root}}^{(\ell,s)}$
& Gewichtsmatrix für den eigenen Knotenzustand \\

$W_{\mathrm{msg}}^{(\ell,s)}$
& Gewichtsmatrix für aggregierte Nachrichten \\

$b^{(\ell,s)}$
& Biasvektor der GNN-Schicht $\ell$ \\

$w_{\mathrm{out}}^s,\ b_{\mathrm{out}}^s$
& Gewichtsvektor und Bias der Ausgabeschicht \\

$w_{\mathrm{skip}}^s,\ b_{\mathrm{skip}}^s$
& Gewichtsvektor und Bias der Skip-Verbindung \\

$\underline{h}_{iq}^{(\ell,s)},\ \overline{h}_{iq}^{(\ell,s)}$
& Schranken der aktivierten Knotenrepräsentation \\

$\underline{v}_{ir}^{(\ell,s)},\ \overline{v}_{ir}^{(\ell,s)}$
& Schranken der Präaktivierung \\

\midrule
\multicolumn{2}{l}{\textbf{Entscheidungs-, Hilfs- und abgeleitete Größen}}\\
\midrule

$x_{iq}$
& Eingabefeature $q$ von Operationsknoten $i$ \\

$h_{ir}^{(\ell,s)}$
& aktivierte Knotenrepräsentation nach Schicht $\ell$ \\

$v_{ir}^{(\ell,s)}$
& Präaktivierung vor der ReLU-Funktion \\

$a_{iq}^{(\ell-1,s)}$
& aggregierte Nachrichtensumme \\

$U_{jik}$
& 1, falls $j$ direkter Maschinenvorgänger von $i$ auf $k$ ist \\

$Z_{jikq}^{(\ell-1)}$
& Linearisierung von $U_{jik}h_{jq}^{(\ell-1,\mathrm{FULL})}$ \\

$\delta_{ir}^{(\ell,s)}$
& Binärvariable der ReLU-Linearisierung \\

$g_{ur}^{s}$
& gepoolte Hidden-Repräsentation von Job $u$ \\

$\bar{x}_{uq}$
& gepooltes Eingabefeature für die Skip-Verbindung \\

$o_u^s$
& ungeclippte affine Ausgabe für Job $u$ \\

$\widehat{P}_{u}^{s}$
& auf $[0,1]$ begrenzte GNN-Schätzung für Job $u$ \\

\bottomrule
\end{tabularx}
\end{table}

Für alle Varianten wird die Eingaberepräsentation durch
$h_{iq}^{(0,s)}=x_{iq}$ initialisiert. Die fünf Eingabefeatures beschreiben
den normierten nominalen Start- und Fertigstellungszeitpunkt, die
Bearbeitungszeit relativ zum Weibull-Skalenparameter, eine normierte
Reparaturrate sowie den normierten Weibull-Formparameter. Da diese Größen von
den Maschinenzuordnungs- und Zeitvariablen des FJSP abhängen, entsteht eine
direkte Verbindung zwischen Ablaufplanung und GNN-Prognose.


% ============================================================
% Architekturabhängige Nachrichtenaggregation
% ============================================================

Die drei untersuchten Architekturvarianten unterscheiden sich ausschließlich
hinsichtlich der Informationen, die zwischen den Operationsknoten
ausgetauscht werden. Für jede Variante wird ein separates neuronales Netz
trainiert und anschließend unabhängig in das Optimierungsmodell eingebettet.
Die nachfolgenden Formulierungen stellen somit alternative und nicht
gleichzeitig verwendete Architekturvarianten dar.


% ============================================================
% Variante 1: Lineare Baseline
% ============================================================

\paragraph{Variante 1: Lineare Baseline ohne Message Passing.}

Die lineare Baseline verarbeitet jeden Operationsknoten unabhängig von seinen
Vorgängern und Nachbarn. Da keine Nachrichten aggregiert werden, hängt die
neue Knotenrepräsentation ausschließlich von der Repräsentation desselben
Knotens in der vorherigen Schicht ab. Die Präaktivierung lautet

\begin{align}
v_{ir}^{(\ell,\mathrm{LIN})}
&=
b_r^{(\ell,\mathrm{LIN})}
+
\sum_{q\in\mathcal{Q}_{\ell-1}}
W_{\mathrm{lin},rq}^{(\ell,\mathrm{LIN})}
h_{iq}^{(\ell-1,\mathrm{LIN})}
\nonumber\\
&&&
\forall\,i\in\mathcal{O},\;
r\in\mathcal{Q}_\ell,\;
\ell\in\mathcal{L}.
\label{eq:gnn-linear-layer}
\end{align}

Diese Variante dient als strukturfreie Baseline. Der Vergleich mit den
nachfolgenden Varianten zeigt, ob die Berücksichtigung der Graphstruktur einen
zusätzlichen Beitrag zur Vorhersagequalität leistet.


% ============================================================
% Variante 2: Jobpräzedenzkanten
% ============================================================

\paragraph{Variante 2: Message Passing über feste Jobpräzedenzkanten.}

Die zweite Variante berücksichtigt die technologisch vorgegebenen
Präzedenzbeziehungen innerhalb eines Jobs. Die Menge $\mathcal{P}_i$ enthält
die direkten Jobvorgänger der Operation $i$. Für jeden Featurekanal $q$
werden deren Repräsentationen durch Summation aggregiert:

\begin{align}
a_{iq}^{(\ell-1,\mathrm{JOB})}
&=
\sum_{j\in\mathcal{P}_i}
h_{jq}^{(\ell-1,\mathrm{JOB})}
\nonumber\\
&&&
\forall\,i\in\mathcal{O},\;
q\in\mathcal{Q}_{\ell-1},\;
\ell\in\mathcal{L}.
\label{eq:gnn-job-aggregation}
\end{align}

Die Präaktivierung eines Knotens setzt sich anschließend aus der eigenen
Repräsentation und den aggregierten Nachrichten seiner Jobvorgänger zusammen:

\begin{align}
v_{ir}^{(\ell,\mathrm{JOB})}
&=
b_r^{(\ell,\mathrm{JOB})}
+
\sum_{q\in\mathcal{Q}_{\ell-1}}
W_{\mathrm{root},rq}^{(\ell,\mathrm{JOB})}
h_{iq}^{(\ell-1,\mathrm{JOB})}
\nonumber\\
&\quad+
\sum_{q\in\mathcal{Q}_{\ell-1}}
W_{\mathrm{msg},rq}^{(\ell,\mathrm{JOB})}
a_{iq}^{(\ell-1,\mathrm{JOB})}
\nonumber\\
&&&
\forall\,i\in\mathcal{O},\;
r\in\mathcal{Q}_\ell,\;
\ell\in\mathcal{L}.
\label{eq:gnn-job-layer}
\end{align}

Da die Jobpräzedenzkanten unabhängig von den Planungsentscheidungen
feststehen, erfordert ihre Aktivierung keine zusätzlichen Binärvariablen. Die
Binärvariablen $\delta_{ir}^{(\ell,\mathrm{JOB})}$ werden ausschließlich für
ReLU-Einheiten benötigt, deren Präaktivierungsintervall die Null enthält.


% ============================================================
% Variante 3: Job- und Maschinenkanten
% ============================================================

\paragraph{Variante 3: Message Passing über Job- und Maschinenkanten.}

Die dritte Variante erweitert die festen Jobpräzedenzkanten um
entscheidungsabhängige Maschinenkanten. Die Binärvariable $U_{jik}$ nimmt den
Wert eins an, wenn die Operationen $j$ und $i$ Maschine $k$ zugeordnet sind
und $j$ dort unmittelbarer Vorgänger von $i$ ist. Eine Nachricht von $j$ nach
$i$ wird folglich nur dann berücksichtigt, wenn die entsprechende
Maschinenkante aktiv ist. Die zur Definition der unmittelbaren Vorgänger
erforderlichen Zuordnungs-, Reihenfolge-, Fluss- und Randbedingungen sind
Bestandteil der Kopplung mit dem FJSP-Grundmodell.

Das Produkt aus der binären Kantenvariable $U_{jik}$ und der kontinuierlichen
Knotenrepräsentation $h_{jq}^{(\ell-1,\mathrm{FULL})}$ wird durch die
Hilfsvariable

\begin{equation}
Z_{jikq}^{(\ell-1)}
=
U_{jik}h_{jq}^{(\ell-1,\mathrm{FULL})}
\label{eq:gnn-full-product-definition}
\end{equation}

repräsentiert. Unter Verwendung der gültigen Schranken
$\underline{h}_{jq}^{(\ell-1,\mathrm{FULL})}$ und
$\overline{h}_{jq}^{(\ell-1,\mathrm{FULL})}$ wird das
Binär-Kontinuierlich-Produkt durch die folgenden vier Ungleichungen exakt
linearisiert:

\begin{align}
Z_{jikq}^{(\ell-1)}
&\ge
\underline{h}_{jq}^{(\ell-1,\mathrm{FULL})}U_{jik}
&&
\forall\,(j,i,k)\in\mathcal{E}^{M},\;
q\in\mathcal{Q}_{\ell-1},\;
\ell\in\mathcal{L},
\label{eq:gnn-full-product-lb1}
\\
Z_{jikq}^{(\ell-1)}
&\le
\overline{h}_{jq}^{(\ell-1,\mathrm{FULL})}U_{jik}
&&
\forall\,(j,i,k)\in\mathcal{E}^{M},\;
q\in\mathcal{Q}_{\ell-1},\;
\ell\in\mathcal{L},
\label{eq:gnn-full-product-ub1}
\\
Z_{jikq}^{(\ell-1)}
&\ge
h_{jq}^{(\ell-1,\mathrm{FULL})}
-
\overline{h}_{jq}^{(\ell-1,\mathrm{FULL})}(1-U_{jik})
&&
\forall\,(j,i,k)\in\mathcal{E}^{M},\;
q\in\mathcal{Q}_{\ell-1},\;
\ell\in\mathcal{L},
\label{eq:gnn-full-product-lb2}
\\
Z_{jikq}^{(\ell-1)}
&\le
h_{jq}^{(\ell-1,\mathrm{FULL})}
-
\underline{h}_{jq}^{(\ell-1,\mathrm{FULL})}(1-U_{jik})
&&
\forall\,(j,i,k)\in\mathcal{E}^{M},\;
q\in\mathcal{Q}_{\ell-1},\;
\ell\in\mathcal{L}.
\label{eq:gnn-full-product-ub2}
\end{align}

Für $U_{jik}=0$ erzwingen die Nebenbedingungen
\eqref{eq:gnn-full-product-lb1}--\eqref{eq:gnn-full-product-ub2} den Wert
$Z_{jikq}^{(\ell-1)}=0$. Für $U_{jik}=1$ folgt dagegen
$Z_{jikq}^{(\ell-1)}=h_{jq}^{(\ell-1,\mathrm{FULL})}$. Damit wird die
Nachricht exakt entsprechend der im Ablaufplan gewählten Maschinenkante ein-
oder ausgeschaltet.

Die vollständige Aggregation umfasst sowohl die Repräsentationen der festen
Jobvorgänger als auch die über $U_{jik}$ aktivierten Nachrichten der
Maschinenvorgänger:

\begin{align}
a_{iq}^{(\ell-1,\mathrm{FULL})}
&=
\sum_{j\in\mathcal{P}_i}
h_{jq}^{(\ell-1,\mathrm{FULL})}
+
\sum_{(j,i,k)\in\mathcal{E}^{M}}
Z_{jikq}^{(\ell-1)}
\nonumber\\
&&&
\forall\,i\in\mathcal{O},\;
q\in\mathcal{Q}_{\ell-1},\;
\ell\in\mathcal{L}.
\label{eq:gnn-full-aggregation}
\end{align}

Die Präaktivierung ergibt sich wiederum aus der eigenen Knotenrepräsentation
und den aggregierten Nachrichten:

\begin{align}
v_{ir}^{(\ell,\mathrm{FULL})}
&=
b_r^{(\ell,\mathrm{FULL})}
+
\sum_{q\in\mathcal{Q}_{\ell-1}}
W_{\mathrm{root},rq}^{(\ell,\mathrm{FULL})}
h_{iq}^{(\ell-1,\mathrm{FULL})}
\nonumber\\
&\quad+
\sum_{q\in\mathcal{Q}_{\ell-1}}
W_{\mathrm{msg},rq}^{(\ell,\mathrm{FULL})}
a_{iq}^{(\ell-1,\mathrm{FULL})}
\nonumber\\
&&&
\forall\,i\in\mathcal{O},\;
r\in\mathcal{Q}_\ell,\;
\ell\in\mathcal{L}.
\label{eq:gnn-full-layer}
\end{align}

Damit bilden die drei Varianten eine schrittweise Erweiterung der verwendeten
Strukturinformationen:

\[
\underbrace{\mathrm{LIN}}_{\text{keine Nachbarschaft}}
\;\longrightarrow\;
\underbrace{\mathrm{JOB}}_{\text{feste Jobpräzedenzkanten}}
\;\longrightarrow\;
\underbrace{\mathrm{FULL}}_{
  \substack{\text{feste Jobpräzedenzkanten}\\
            +\text{ entscheidungsabhängige Maschinenkanten}}
}.
\]

\paragraph{Linearisierung der ReLU-Aktivierung.}

Auf die jeweilige Präaktivierung $v_{ir}^{(\ell,s)}$ wird in allen drei
Architekturvarianten dieselbe ReLU-Funktion angewendet:

\begin{equation}
h_{ir}^{(\ell,s)}
=
\max\{0,v_{ir}^{(\ell,s)}\}.
\label{eq:gnn-relu-definition}
\end{equation}

Falls
$\underline{v}_{ir}^{(\ell,s)}<0<\overline{v}_{ir}^{(\ell,s)}$ gilt, wird
diese stückweise lineare Funktion mithilfe der Binärvariable
$\delta_{ir}^{(\ell,s)}$ exakt abgebildet:

\begin{align}
h_{ir}^{(\ell,s)}
&\ge v_{ir}^{(\ell,s)},
\label{eq:gnn-relu-lb-pre}\\
h_{ir}^{(\ell,s)}
&\ge 0,
\label{eq:gnn-relu-lb-zero}\\
h_{ir}^{(\ell,s)}
&\le
v_{ir}^{(\ell,s)}
-\underline{v}_{ir}^{(\ell,s)}
\bigl(1-\delta_{ir}^{(\ell,s)}\bigr),
\label{eq:gnn-relu-ub-active}\\
h_{ir}^{(\ell,s)}
&\le
\overline{v}_{ir}^{(\ell,s)}\delta_{ir}^{(\ell,s)},
\label{eq:gnn-relu-ub-inactive}\\
\delta_{ir}^{(\ell,s)}
&\in\{0,1\}.
\label{eq:gnn-relu-binary}
\end{align}

Die Indizes $i\in\mathcal{O}$, $r\in\mathcal{Q}_{\ell}$,
$\ell\in\mathcal{L}$ und $s\in\mathcal{A}$ gelten für die
Nebenbedingungen \eqref{eq:gnn-relu-lb-pre}--\eqref{eq:gnn-relu-binary}.
Für $\delta_{ir}^{(\ell,s)}=0$ folgt $h_{ir}^{(\ell,s)}=0$; für
$\delta_{ir}^{(\ell,s)}=1$ gilt
$h_{ir}^{(\ell,s)}=v_{ir}^{(\ell,s)}$. Liegt das gesamte
Präaktivierungsintervall auf einer Seite der Null, kann die Binärvariable
entfallen: Für $\overline{v}_{ir}^{(\ell,s)}\le 0$ wird
$h_{ir}^{(\ell,s)}=0$ fixiert, während für
$\underline{v}_{ir}^{(\ell,s)}\ge 0$ die lineare Gleichung
$h_{ir}^{(\ell,s)}=v_{ir}^{(\ell,s)}$ ausreicht. Enge, neuronenspezifische
Schranken stärken somit nicht nur die LP-Relaxation, sondern können auch die
Anzahl zusätzlicher Binärvariablen reduzieren.

\paragraph{Jobspezifisches Pooling und Ausgabeschicht.}

Nach der letzten GNN-Schicht werden die Knotenrepräsentationen aller
Operationen eines Jobs addiert. Für jeden Job $u$ und Hidden-Kanal $r$ gilt

\begin{align}
g_{ur}^{s}
&=
\sum_{i\in\mathcal{O}_u}h_{ir}^{(L,s)},
&&
r\in\mathcal{Q}_{L},\quad s\in\mathcal{A}.
\label{eq:gnn-job-pooling}
\end{align}

Parallel dazu werden die ursprünglichen Eingabefeatures für die direkte
Skip-Verbindung jobspezifisch aggregiert:

\begin{align}
\bar{x}_{uq}
&=
\sum_{i\in\mathcal{O}_u}x_{iq},
&&
q\in\mathcal{Q}_{0}.
\label{eq:gnn-input-pooling}
\end{align}

Die skalare Ausgabe des Netzes ergibt sich aus der gepoolten
Hidden-Repräsentation und der gepoolten Eingaberepräsentation:

\begin{align}
o_u^s
&=
b_{\mathrm{out}}^s+b_{\mathrm{skip}}^s
+\sum_{r\in\mathcal{Q}_{L}}w_{\mathrm{out},r}^s g_{ur}^{s}
+\sum_{q\in\mathcal{Q}_{0}}w_{\mathrm{skip},q}^s\bar{x}_{uq}.
\label{eq:gnn-output}
\end{align}

Im trainierten Netz wird diese rohe Ausgabe zunächst durch eine ReLU auf
nichtnegative Werte und anschließend durch Clipping auf eins begrenzt:

\begin{equation}
\widehat{P}_{u}^{s}
=
\min\bigl\{1,\max\{0,o_u^s\}\bigr\}.
\label{eq:gnn-probability-postprocessing}
\end{equation}

Für ein vorgegebenes Serviceniveau $\alpha_u\in(0,1]$ ist die Bedingung
$\widehat{P}_{u}^{s}\ge\alpha_u$ äquivalent zu

\begin{equation}
o_u^s\ge\alpha_u
\qquad
\forall\,u.
\label{eq:gnn-service-level}
\end{equation}

Die ReLU- und Clipping-Operation der Ausgabeschicht müssen daher nicht
zusätzlich in das MILP aufgenommen werden. Dies vermeidet weitere
Binärvariablen, ohne die zulässige Menge bezüglich der
Serviceniveauanforderung zu verändern. Die auf $[0,1]$ begrenzte Schätzung
aus \eqref{eq:gnn-probability-postprocessing} wird lediglich zur Interpretation
und Ausgabe der Prognose benötigt. Da das GNN auf simulierten
Pünktlichkeitsanteilen trainiert wird, ist $\widehat{P}_{u}^{s}$ eine
statistische Punktschätzung und keine mathematisch garantierte Untergrenze der
tatsächlichen Pünktlichkeitswahrscheinlichkeit. Die Güte der resultierenden
Serviceniveaubedingung ist daher empirisch anhand unabhängiger
Simulationsdaten zu validieren.

\paragraph{Einbindung in das FJSP-Grundmodell.}

Die Nebenbedingungen des deterministischen FJSP-Grundmodells zur eindeutigen
Maschinenzuordnung, zu den Bearbeitungs- und Fertigstellungszeiten, zu den
Jobpräzedenzen sowie zum Ausschluss zeitlicher Überlappungen bleiben
unverändert bestehen. Für die FULL-Variante werden zusätzlich die Variablen
$U_{jik}$ und die zugehörigen Kopplungsbedingungen eingeführt, damit $U_{jik}$
genau die im Ablaufplan aktiven unmittelbaren Maschinenvorgänger beschreibt.
Die Gleichungen der jeweils gewählten GNN-Variante, die
ReLU-Linearisierungen, das jobspezifische Pooling und die
Serviceniveaubedingung \eqref{eq:gnn-service-level} erweitern dieses
Grundmodell.

Die nichtlinearen Nebenbedingungen des analytischen Referenzmodells zur
Berechnung der Pünktlichkeitswahrscheinlichkeiten werden im GNN-Modell nicht
beibehalten. Ihre Funktion übernimmt das trainierte Ersatzmodell. Da nach dem
Training sämtliche Netzwerkparameter konstant sind und alle verbleibenden
Nichtlinearitäten exakt linearisiert beziehungsweise für die
Serviceniveaubedingung äquivalent umgeformt werden, ist das resultierende
Optimierungsmodell ein MILP.
